\section{Colour Key and Style Guide}
\label{app:colors}

This section explains the visual conventions used throughout the survey for boxes, emphasis, and structural elements.

\subsection{Box Types and Colours}

\begin{definitionbox}[Definition Box (Dark Blue)]
\textbf{Definition boxes} introduce core concepts and provide precise technical definitions. Use these when a term needs clear specification before use.\\

Example concepts: ``Feature'', ``Circuit'', ``Sparse Autoencoder''.\\

Style: Dark blue border, light blue background.
\end{definitionbox}

\begin{examplebox}[Example Box (Green)]
\textbf{Example boxes} provide concrete illustrations, worked examples, or case studies. Use these to make abstract concepts concrete.\\

Example contents: ``Induction in a Repeated Sequence'', ``Example Application of Steering Vectors'', ``Probing Classifiers for Sleeper Agents''.\\

Style: Green border, light green background.
\end{examplebox}

\begin{insightbox}[Insight Box (Gold)]
\textbf{Insight boxes} explain intuitions behind mathematical or conceptual results. Use these to build understanding of why something works, not just that it works.\\

Example contents: ``Trade-offs Between Variants'', ``The Complexity--Interpretability Trade-off''.\\

Style: Gold border, light gold background.
\end{insightbox}

\begin{warningbox}[Warning Box (Red)]
\textbf{Warning boxes} highlight critical limitations, failure modes, or common misconceptions. Use these to prevent readers from adopting incorrect or unsafe interpretations.\\

Example contents: ``Caveats on CLT Fidelity Claims'', ``The Probe Causality Problem'', ``The Partial Coverage Issue''.\\

Style: Red border, light red background.
\end{warningbox}

\begin{practicebox}[Practice Box (Teal)]
\textbf{Practice boxes} illustrate suggestions that are actionable. These include protocol recommendations, checklists for verifying approaches, and step-by-step decision-making walkthroughs. Use these for practical implementations of procedures.\\

Example contents: ``Guardrail Framework for Automated Interpretability'', ``Validation Procedures for Introspective Models'', ``Recommended Evaluation Protocol''.\\

Style: Teal border, light teal background.
\end{practicebox}

\subsection{Citation Evidence Levels}
\label{app:citation-colors}

This survey cites both peer-reviewed publications and non-peer-reviewed research outputs. To help readers assess the maturity and review status of the evidence, citations are colour-coded:

\begin{itemize}
    \item \textcolor{citeestablished}{\textbf{Orange citations}} denote formally
    published work: peer-reviewed conference and journal papers, and books.

    \item \textcolor{citeemerging}{\textbf{Green citations}}: Research outputs that have not undergone formal peer review, including preprints, technical reports, workshop papers, blog posts, and other emerging resources.
\end{itemize}

Green citations should not automatically be interpreted as lower quality or less reliable. Some originate from leading research organisations and may have undergone substantial internal review. However, they have not been subjected to the formal academic peer-review process and should therefore be interpreted with appropriate caution. Readers are encouraged to verify the current publication status of cited work where relevant.


